\begin{definition}[Characteristic convention]\label{mod:basic-characteristicconvention}
Let \(E\) be an additive monoid with one and let \(p\in\mathbf N\).
If \(E\) has characteristic \(p\), then the canonical image of \(p\)
in \(E\) is zero:
\[
 (p:E)=0.
\]
In particular, in characteristic two,
\[
 (2:E)=0.
\]
If \(E\) has characteristic zero, every nonzero natural number has
nonzero canonical image:
\[
 n\ne0\quad\Longrightarrow\quad(n:E)\ne0.
\]
When \(E\) is an additive group with one, the analogous assertion holds
for the canonical integer map \(\mathbf Z\to E\):
\[
 z\ne0\quad\Longrightarrow\quad(z:E)\ne0.
\]
These statements concern only canonical casts.  They do not identify
\(p\) with the characteristic of the residue field of a local field.
Later equal-characteristic branches supply \texttt{CharP E p}, while
mixed-characteristic branches supply \texttt{CharZero E} together with
their local-field hypotheses.
\LeanFile{LanglandsFirstMainLemma/Basic/CharacteristicConvention.lean}
\lean{LanglandsFirstMainLemma.characteristicConvention}
\leanok
\SourceText{Remark rem:characteristic-convention and Introduction, subsection Equal characteristic.}
\uses{mod:prelude}
\FormalizationContract
\end{definition}
